\section{Set-Covering Formulations}\label{sec:formulations}

Let $\mathcal R$ be the set of all feasible elementary internal sequences. For
$r\in\mathcal R$, let $a_{jr}$ be 1 if sequence $r$ contains job $j$ and 0
otherwise, and let $x_r$ indicate whether the sequence is selected. Both
formulations use the machine-capacity constraint
\begin{equation}
    \sum_{r\in\mathcal R}x_r\le m.
    \label{eq:machine-capacity}
\end{equation}

\subsection{Formulation with explicit outsourcing variables}
\label{subsec:master-outsourcing-formulation}

The default formulation generates only internal schedule columns. Let $y_j$ be 1
if job $j$ is outsourced. The nonlinear set-covering formulation is
\begin{align}
(\mathrm{SC\text{-}M})\qquad
\min\quad
    & \sum_{r\in\mathcal R}c_r^{\mathrm{in}}x_r
      +G\left(\sum_{j\in\mathcal J}b_jy_j\right)
      \label{eq:scm-obj}\\
\text{s.t.}\quad
    & \sum_{r\in\mathcal R}a_{jr}x_r+y_j\ge1
      && j\in\mathcal J, \label{eq:scm-cover}\\
    & \sum_{r\in\mathcal R}x_r\le m, \label{eq:scm-machines}\\
    & x_r\in\{0,1\}
      && r\in\mathcal R, \label{eq:scm-x}\\
    & y_j\in\{0,1\}
      && j\in\mathcal J. \label{eq:scm-y}
\end{align}

The implemented master represents $G$ by tariff-segment variables. Let $z_\ell$
select segment $\ell$ and let $q_\ell$ be the baseline outsourcing amount assigned
to that segment. The tariff term in \eqref{eq:scm-obj} is replaced by
\begin{equation}
    \sum_{\ell\in\mathcal L}
    \left(a_\ell q_\ell+c_\ell z_\ell\right),
    \label{eq:tariff-objective}
\end{equation}
and the following constraints are added:
\begin{align}
    &\sum_{j\in\mathcal J}b_jy_j
      =\sum_{\ell\in\mathcal L}q_\ell,
      \label{eq:tariff-balance}\\
    &\underline q_\ell z_\ell
      \le q_\ell
      \le\overline q_\ell z_\ell
      &&\ell\in\mathcal L,
      \label{eq:tariff-bounds}\\
    &\sum_{\ell\in\mathcal L}z_\ell=1,
      \label{eq:tariff-one-segment}\\
    &z_\ell\in\{0,1\},\quad q_\ell\ge0
      &&\ell\in\mathcal L.
      \label{eq:tariff-domains}
\end{align}
Only the internal sequence variables are generated dynamically. The outsourcing and
tariff decisions remain in the restricted master problem.

Let $\pi_j\ge0$ be the dual value of covering row
\eqref{eq:scm-cover}, and let $\sigma\le0$ be the dual value of
\eqref{eq:scm-machines} under the displayed minimization convention. The reduced
cost of an internal schedule $r$ is
\begin{equation}
    \bar c_r^{\mathrm{in}}
    =
    c_r^{\mathrm{in}}
    -\sum_{j\in\mathcal J}\pi_j a_{jr}
    -\sigma.
    \label{eq:internal-reduced-cost}
\end{equation}
Branching rows and active cuts add the corresponding dual terms; these terms are
described with the algorithm in Section~\ref{sec:solution}.

\subsection{Formulation with outsourcing columns}
\label{subsec:columnized-outsourcing-formulation}

The alternative formulation represents an outsourced set by one column. Let
$\mathcal O=2^{\mathcal J}\setminus\{\varnothing\}$, let $d_{jo}$ indicate whether
job $j$ belongs to outsourced set $o$, and define
\begin{equation}
    c_o^{\mathrm{out}}
    =
    G\left(\sum_{j\in\mathcal J}b_jd_{jo}\right).
\end{equation}
Variable $w_o$ selects outsourcing column $o$. The formulation is
\begin{align}
(\mathrm{SC\text{-}C})\qquad
\min\quad
    &\sum_{r\in\mathcal R}c_r^{\mathrm{in}}x_r
     +\sum_{o\in\mathcal O}c_o^{\mathrm{out}}w_o
     \label{eq:scc-obj}\\
\text{s.t.}\quad
    &\sum_{r\in\mathcal R}a_{jr}x_r
     +\sum_{o\in\mathcal O}d_{jo}w_o\ge1
     &&j\in\mathcal J, \label{eq:scc-cover}\\
    &\sum_{r\in\mathcal R}x_r\le m,
     \label{eq:scc-machines}\\
    &\sum_{o\in\mathcal O}w_o\le1,
     \label{eq:scc-at-most-one-outsourcing}\\
    &x_r\in\{0,1\}
     &&r\in\mathcal R,\\
    &w_o\in\{0,1\}
     &&o\in\mathcal O.
\end{align}
No empty outsourcing column is needed: selecting no $w_o$ represents the case in
which all jobs are processed internally.

Let $\mu\le0$ be the dual value of
\eqref{eq:scc-at-most-one-outsourcing}. The reduced cost of outsourcing column $o$
is
\begin{equation}
    \bar c_o^{\mathrm{out}}
    =
    G\left(\sum_{j\in\mathcal J}b_jd_{jo}\right)
    -\sum_{j\in\mathcal J}\pi_jd_{jo}
    -\mu.
    \label{eq:outsourcing-reduced-cost}
\end{equation}
Thus, after removing the constant $-\mu$, outsourcing pricing solves
\begin{equation}
    \max_{O\subseteq\mathcal J}
    \left\{
        \sum_{j\in O}\pi_j
        -
        G\left(\sum_{j\in O}b_j\right)
    \right\}.
    \label{eq:outsourcing-subset-pricing}
\end{equation}

\subsection{Comparison of the formulations}

Formulation $\mathrm{SC\text{-}M}$ has one pricing family and a larger master that
contains the tariff representation. Formulation $\mathrm{SC\text{-}C}$ keeps the
master outsourcing structure compact but requires a second pricing problem and
membership branching for outsourced sets. The two formulations therefore expose
different computational trade-offs: master complexity and tariff branching in
$\mathrm{SC\text{-}M}$, versus an additional combinatorial pricing family in
$\mathrm{SC\text{-}C}$. Their internal schedule columns and internal pricing problem
are identical.
